\section{问题二模型建立与求解}
\label{sec:question-two}

问题一把每个子图看作独立Task，子图边界可能带来写回、读入和激活等待。问题二仍提交操作归属、分核和核内子图序列，但同一核心上的子图合成一个核级Task。数据可在该核心的L1或UB内继续使用，跨核心的数据仍经DDR传送。因此，本问的决策不只是减少跨核边：把消费者集中在同一核心可以节省读入，也可能让几个大张量同时留在片上，随后发生换出与恢复。本节依次建立核心级结构搬运、物理副本容量、有限候选求解和逐图结果之间的关系。

\subsection{数据预处理与核心驻留域}
\subsubsection{核心驻留域划分}
给定第\ref{sec:question-one}节的计划$X=(g,a,\sigma)$，先由\texttt{node\_to\_subgraph}取得操作所属子图$g(v)$，再从\texttt{core\_schedules}取得子图所在核心$a(g(v))$及同核先后次序。原计算图中的COPY操作不由外层计划重新分配；候选只覆盖非COPY计算操作，实际输入、输出和跨核COPY由评估程序根据方案生成。每份候选须通过式\eqref{eq:plan-coverage}和式\eqref{eq:plan-sequence-unique}的覆盖与唯一性检查，核级Task、物理版本和事件时间线再由评估程序展开。对每个逻辑张量$t$，索引保存其大小$b_t$、生产者$p(t)$、去重后的计算消费者$U(t)$、存储区位置以及最终写回指示$o_t$。本研究使用的100张图均按同一字段定义建立这些关系，候选的归属与顺序变化不改变原计算依赖。

问题一按消费Task去重，问题二按消费核心去重。定义计算操作$v$的驻留域$d_B(v)=a(g(v))$，则张量$t$的核级消费域为
\begin{equation}
\mathcal D_t^B=\{d_B(v):v\in U(t)\}.
\label{eq:b-domain}
\end{equation}
同一张量在一个核心上的多个计算消费者，在式\eqref{eq:b-domain}中只占一个元素。这是接收核心去重，不表示张量跨不同核心共享同一片上物理副本。一个核心上的子图即使被分成多个块，其生产与消费边仍保留；评估程序的核内排序按提交的子图顺序安排这些操作。张量的逻辑编号用于统计接收域，容量检查则使用换出、恢复形成的物理版本编号。

\begin{table}[H]\centering\caption{问题一与问题二的驻留及同步边界}\label{tab:b-rule-contrast}
\begin{tabularx}{\linewidth}{lYY}\toprule
对象 & 问题一：Task驻留 & 问题二：核心驻留\\\midrule
片上副本可复用范围 & 单个子图Task & 同核合并后的Task\\
相同输入的读入次数 & 按消费子图去重 & 按消费核心去重\\
中间结果的跨域COPY & 源Task写出、各远端Task读入 & 每个远端接收核心各有写出和读入一对\\
跨核释放起点 & 前驱Task完成并等待1000 cycle & 源端COPY\_OUT完成并等待500 cycle\\
容量统计 & Task内部的物理副本 & 同核合并序上的物理副本\\\bottomrule
\end{tabularx}\end{table}

\subsubsection{跨核结构搬运}
问题二允许同核多个子图复用片上数据，因而同一个张量在一个核心上只需一个驻留副本。对输入$t$记$I_B(t)=\{a(g(v)):v\in U(t)\}$；对有生产者的$t$记$R_B(t)=\{a(g(v)):v\in U(t),\,a(g(v))\ne a(g(p(t)))\}$。每个远端接收核需要源核写出与目标核读入，最终输出另行写回，得到
\begin{equation}
Q_B^0=Q_C^0=
\sum_{t\in\mathcal T_{\rm in}}b_t|I_B(t)|
+\sum_{t\in\mathcal T_{\rm prod}}b_t(2|R_B(t)|+o_t).
\label{eq:b-bytes}
\end{equation}
式\eqref{eq:b-bytes}中，图输入没有计算生产者，每个接收核心从DDR读入一次，贡献$b_t|I_B(t)|$。中间张量在生产核心形成，每个远端接收核心均需一份源写出与一份目标读入，贡献$2b_t|R_B(t)|$；原图要求输出的张量另写回一次，贡献$b_t o_t$。等式$Q_B^0=Q_C^0$指同一计划在无Cache与有Cache场景下的结构COPY字节相同；问题三改变读服务来源和完成时刻，不删除这条逻辑搬运。

若一份$b$字节输入由同核两个子图消费，问题一可能读入$2b$，问题二只读入$b$。若一份中间结果的生产者位于核心0、两个消费子图分别位于核心1和2，本问需要两对跨核COPY，即$4b$；将两个消费者放到同一个远端核心后变为$2b$，全部放在生产核心则该项为零。问题一对两个远端Task的同一张量只需源Task写出一次、两个Task各读入一次，即$3b$。这些比较只涉及结构搬运，集中数据后可能产生的Spill另行计入。对于既要跨核发送又要最终写回的张量，评估程序分别创建跨核COPY与最终输出COPY，因此式中保留$o_t$。

上述例子说明$Q_B^0$与$Q_A^0$没有固定的大小关系。令同一生产张量$t$的外部消费Task数为$r_A=|R_A(t)|$，远端消费核心数为$r_B=|R_B(t)|$，并令$o_t$表示最终写回指示。由式\eqref{eq:a-bytes}和式\eqref{eq:b-bytes}，该张量在两种场景下的结构搬运之差为
\begin{equation}
Q_{B,t}^0-Q_{A,t}^0
=b_t\left(2r_B+o_t-r_A-\mathbf 1\{r_A>0\ \text{或}\ o_t=1\}\right).
\label{eq:ab-structure-difference}
\end{equation}
其中$r_A$按Task计数，包含与生产Task同核但属于不同Task的消费者；$r_B$按远端核心计数且不包含生产核心。右端可为正、负或零，且只比较同一计划下的结构COPY，未含Spill、带宽排队和执行时序。因此A/B方案须分别按各自的搬运公式和事件模拟结果评价。

同核切块可以改变执行顺序而不改变式\eqref{eq:b-bytes}的结构字节数，但张量驻留时间和容量恢复仍会变化。跨核目标COPY\_IN需等待源COPY\_OUT完成，再加500 cycle释放延迟。

若$e^{\rm out}_{t,i\to j}$和$s^{\rm in}_{t,i\to j}$分别是该边源COPY\_OUT的结束和目标COPY\_IN的开始时刻，则
\begin{equation}
s^{\rm in}_{t,i\to j}\ge e^{\rm out}_{t,i\to j}+500.
\label{eq:b-copy-release}
\end{equation}
源端写出受MTE3流水线及共享DDR服务影响，其完成时间移动时，目标读入的可发射时间随之移动。

公式中的去重以消费核心为单位。把同一输入的消费者集中到一核可减少重复读入；把中间张量的消费者移回生产核可减少写读一对COPY。两种操作也可能使一个核心的计算和存活张量过于集中，结构搬运下降与最终工期的关系由完整事件时间线决定。

\subsection{物理副本、容量与Spill}
\subsubsection{物理副本存活期}
Buffets把片上缓冲的填充、读取和释放交给显式的编排逻辑管理\cite{pellauer2019buffets}；本题的L1、UB同样按“申请时占用、最后读者完成时释放”管理空间，因此可以用物理副本的存活区间描述容量。核心$k$接收或生产一份逻辑张量后，该副本从片上空间申请时刻$A_{\tilde t}$开始占用L1或UB，直到最后一个读者完成的时刻$L_{\tilde t}$释放。读者包括计算操作以及将结果写入DDR的COPY。若Step2将它换出并在下次使用前恢复，恢复后占用空间的是新的物理版本$\tilde t'$。对存储区$q$，执行过程必须满足
\begin{equation}
\mathrm{Live}_{kq}(\tau)
=\sum_{\tilde t\in\mathcal V_{kq}(X)}b_{\tilde t}
\mathbf1\{A_{\tilde t}\le\tau<L_{\tilde t}\}
\le C_q,\qquad q\in\{\mathrm{L1},\mathrm{UB}\}.
\label{eq:b-physical-live}
\end{equation}
这里$\mathcal V_{kq}(X)$是计划$X$在核心$k$、存储区$q$产生的物理版本集合，容量$C_q$沿用式\eqref{eq:capacity}。分核与核内顺序改变操作的申请和末次使用次序，因而即使式\eqref{eq:b-bytes}的结构字节相同，式\eqref{eq:b-physical-live}中的重叠区间也会不同。

\begin{figure}[H]\centering
\begin{tikzpicture}[x=0.95cm,y=0.95cm]
\draw[alink] (0,0)--(12.0,0) node[right]{$\tau$};
\draw[very thick,blue!70!black] (1.0,0.9)--(10.3,0.9);
\draw[very thick,orange!80!black] (3.1,1.55)--(8.4,1.55);
\foreach \x in {1.0,3.1,8.4,10.3}
  \draw[black!65] (\x,-0.12)--(\x,0.13);
\node[font=\small,below,align=center] at (1.0,-0.16){第一份\\申请};
\node[font=\small,below,align=center] at (3.1,-0.16){第二份\\申请};
\node[font=\small,below,align=center] at (8.4,-0.16){第二份\\末次读};
\node[font=\small,below,align=center] at (10.3,-0.16){第一份\\末次读};
\node[left,font=\small] at (1.0,0.9){$\tilde t_1$};
\node[left,font=\small] at (3.1,1.55){$\tilde t_2$};
\draw[densely dashed,black!50] (3.1,0.2)--(3.1,1.85);
\draw[densely dashed,black!50] (8.4,0.2)--(8.4,1.85);
\node[font=\small] at (5.8,2.05){两份副本同时占用片上空间};
\end{tikzpicture}
\caption{同核两份物理副本的生存区间及容量重叠}\label{fig:b-lifetime}
\end{figure}

例如，两份各81920 B的UB张量若在图\ref{fig:b-lifetime}所示的区间内同时存活，峰值为163840 B，超过131072 B的UB容量32768 B。前移其中一份的末次使用或推迟另一份的申请可以消除这一重叠；否则Step2须通过换出、恢复组织空间。

\subsubsection{换出恢复与容量约束}
ZigZag在探索映射时同时评估存储层次上的数据驻留\cite{mei2021zigzag}。本文也在候选排序中加入容量压力项，但只作粗略估计：完整构图之前，程序在每核的局部拓扑操作序上建立位置索引；每个张量在该核上的局部生产者与消费者中，最早操作位置为代理生存起点，最晚操作位置为代理终点。对计划$X$沿这些位置统计张量的理想存活字节$\widehat{\mathrm{Live}}_{kq}(j)$，以
\begin{equation}
\Pi(X)=\sum_{k,q}\left[\max_j\widehat{\mathrm{Live}}_{kq}(j)-C_q\right]_+
\label{eq:pressure}
\end{equation}
描述未换出时的容量压力，其中$[z]_+=\max\{z,0\}$。该代理对一份张量在某核心只设一段区间，只用于排列候选；实际容量约束由式\eqref{eq:b-physical-live}和评估程序的展开结果确定。评估程序的核内调度依次进行Step1初排及子图优先级调整、Step2插入Spill、Step3补足内存复用依赖；任何一步出现依赖环，候选不进入成功集合。

一次Spill记录的张量大小为$b_s$；恢复COPY\_IN一定搬运$b_s$，若本次换出还需把当前数据写回DDR，则另有一份$b_s$的COPY\_OUT。以输出字段$z_s\in\{0,1\}$表示这次换出是否实际复制数据，Spill字节可从记录逐项复算为
\begin{equation}
Q_B^{\rm spill}=\sum_{s\in\mathcal S_{\rm spill}} b_s(1+z_s).
\label{eq:b-spill}
\end{equation}
相同的容量超额$\Pi$可能因恢复次数、字节数和读者次序不同而产生不同的Spill搬运；Spill还可能推迟关键计算的发射时刻，故工期取决于它在时间线中的位置。

由结构COPY、Spill和原图必要COPY的关系，评估程序输出的调度COPY与额外搬运满足
\begin{equation}
Q_B^{\rm sched}=Q_B^0+Q_B^{\rm spill},\qquad
D_B=Q_B^{\rm sched}-Q^{\rm orig}
=(Q_B^0-Q^{\rm orig})+Q_B^{\rm spill}.
\label{eq:b-breakdown}
\end{equation}
第一项是分核与核级边界产生的结构附加搬运，第二项是容量恢复搬运。字节量以B计，换算MiB时除以$2^{20}$。

表\ref{tab:b-copy-ledger}按式\eqref{eq:b-breakdown}把每个核数下100张图的继承方案逐图分解后取算术平均。额外搬运合计等于结构新增与Spill新增之和，再加原图COPY即为调度COPY总量；该恒等式在500个B结果中逐行成立。

\begin{table}[H]\centering\small
\caption{问题二继承方案的COPY分解（每核数100图；字节列为逐图均值，MiB）}\label{tab:b-copy-ledger}
\begin{tabular}{crrrrr}\toprule
核数 & 原图COPY & 结构新增 & Spill新增 & 额外搬运合计 & 有Spill图数\\\midrule
1 & 2.30 & 0.00 & 24.86 & 24.86 & 81\\
2 & 2.30 & 1.60 & 11.64 & 13.24 & 45\\
3 & 2.30 & 2.50 & 11.39 & 13.89 & 41\\
4 & 2.30 & 3.52 & 10.83 & 14.35 & 38\\
5 & 2.30 & 4.59 & 10.28 & 14.87 & 38\\\bottomrule
\end{tabular}\end{table}

由单核到五核，结构新增均值从0升至4.59 MiB，Spill新增均值从24.86降至10.28 MiB；额外搬运合计在二核时降至13.24 MiB，此后逐核上升至五核的14.87 MiB。单核有81张图发生Spill，五核降至38张：分核缩短了单个核级Task内物理副本的重叠，但也引入了跨核写读。

\subsection{初始方案与求解}
\subsubsection{核级驻留方案}
Kwon等在多DNN加速器上把资源划分与层调度联合优化\cite{kwon2021heterogeneous}；问题二的分核与核内顺序同样相互牵制：分核决定接收域，核内顺序决定物理副本的重叠。整图B0保留原计算图在一个核级Task内的完整驻留，但多核配置中其余核心空闲。B1先求计算依赖的弱连通分量，按分量$\max(M,V)$工作量降序排列，再装入当前桶$\max(M,V)$工作量最小的核心；每个非空核心形成一个子图。无跨分量计算依赖，B1把天然并行的部分分散到不同核心，却无法切开一个占主导的大分量。B2和B3分别使用两种全图就绪序，按剩余$\max(M,V)$工作量目标切成至多$4K$、$8K$个连续块；接近目标的切点优先选择跨切点张量字节较少者。每个新增块接入核心时，对临时计划调用轻量评分，以共同考虑路径、核心负载、结构COPY和容量压力。问题二的块数上限（$4K$、$8K$）比问题一（$2K$、$4K$）更大，因为核级驻留允许同核多块共用片上副本，切得更细不会像问题一那样同比放大结构搬运，但代价是可能加重核内容量压力；更多的切点也给贪心分核更多调整空间。

\begin{table}[H]\centering\caption{问题二四个起点的实际构造及作用}\label{tab:b-seeds}
\begin{tabularx}{\linewidth}{cYY}\toprule
起点 & 构造规则 & 对核级驻留的影响\\\midrule
B0 & 整图一块，放在核心0 & 无新增跨核边界，保留最长驻留区间\\
B1 & 弱连通分量降序装入最轻核心 & 利用独立分量，子图数量不超过非空核数\\
B2 & 沿H序按$\max(M,V)$目标切至多$4K$块 & 优先考虑长计算链，并在候选切点中压低跨切点字节\\
B3 & 沿R序按$\max(M,V)$目标切至多$8K$块 & 将末次消费释放与新输出规模纳入切点选择\\\bottomrule
\end{tabularx}\end{table}

H序使用式\eqref{eq:a-height}的后继最长计算路径优先级；它在每一步仅从已经就绪的操作中选择。R序在此基础上加入逻辑张量的剩余消费者数。令$n_j(t)$为排入第$j$个操作前张量$t$尚未安排的消费者数，从就绪集合中先取H优先级最高的至多8个，组成$\mathcal R_j^{(8)}$。然后计算
\begin{equation}
r_j(v)=\sum_{\substack{t:\,v\in U(t)\\n_j(t)=1}}b_t
-\sum_{t:\,p(t)=v}b_t,
\qquad v_j^*=\mathop{\arg\max}_{v\in\mathcal R_j^{(8)}}
\bigl(r_j(v),h(v),-\operatorname{id}(v)\bigr).
\label{eq:b-r-order}
\end{equation}
第一项只统计本步成为末次消费者的输入，第二项是本步新产生的输出大小。若多个操作释放量相同，先选H优先级高者，再选编号小者。R序是全图静态构造顺序，物理副本的实际生存区间仍由评估程序展开确定。

\subsubsection{方案调整与候选筛选}
每个起点保留原方案，并尝试迁块、同核合并、调序和容量风险块的前后移动。问题一同核数的最终方案也作为问题二的热启动锚点，用于直接比较场景变化带来的工期差异。迁块从累计计算周期最大的核心取其列表首块，移至其他核心中最轻者的列表末端；合并保持同核驻留域，却改变提交边界；调序不改变操作归属与式\eqref{eq:b-bytes}，但可改变式\eqref{eq:b-physical-live}中的重叠。每份编辑方案先规范化，再检查操作覆盖、商图无环及同核直接依赖，非法方案跳过；跨起点产生的相同计划按内容去重。

对合法候选，记$L_k(X)$为核心$k$的计算周期总和，$\widehat L_B(X)$为沿原计算拓扑序的路径代理，跨核依赖在代理中附加500 cycle。令$\widehat Q_B(X)$为候选求解器的静态COPY代理，$\Pi(X)$为式\eqref{eq:pressure}，则入围排序采用
\begin{equation}
\widehat F_B(X)=\max\left\{\widehat L_B(X),\max_k L_k(X),
\left\lceil\frac{\widehat Q_B(X)}{60}\right\rceil\right\},\qquad
\widehat\kappa_B(X)=\bigl(\widehat F_B,\widehat Q_B,\Pi\bigr).
\label{eq:b-light-score}
\end{equation}
$\widehat Q_B$按逻辑张量及其接收核心快速累加，$\widehat L_B$未模拟源端COPY耗时、流水线争用和共享DDR，因此三项只用于候选排序，最终计划按评估程序返回的$F_B$和$D_B$选择。

\subsubsection{完整评价与方案选定}
完整评价先覆盖去重后的整图、分量装桶、上一核已选计划和问题一热启动计划，再从其余候选中至少评价一个代表。设去重后必评计划数为$m$、请求额度为$q$，常规完整调用额度为$\max\{\min(q,2),m+1\}$。评估程序的展开阶段构造核级Task及结构COPY，执行核内Step1/2/3并返回COPY与Spill统计；展开失败时保留错误文本并尝试下一份候选，不消耗完整模拟额度。展开成功后才调用B场景事件模拟，计算每条流水线、DDR及500 cycle跨核释放共同决定的$F_B$。完整模拟即使失败也占一次调用额度；若当前最优慢于上一核数的继承计划，则补评尚未评价的对应回退计划。统一单核基准的回退补评仅在A场景执行。最终在全部成功结果中按$(F_B,D_B)$字典序取方案；正文逐图B结果再按$1,\ldots,K$核计划与整图基准取最短工期并补齐空核。

\begin{algorithm}[H]\caption{问题二的核级驻留候选与完整评价}\label{alg:b-solve}
\begin{algorithmic}[1]
\REQUIRE 原始图$G$、核心数$K$、题设容量和带宽参数及请求完整调用额度$q$
\ENSURE 最终计划$X_B^*$、工期$F_B$、结构及Spill搬运
\STATE 从原图建立计算依赖、张量生产消费索引、H序与R序
\STATE 生成B0整图、B1分量装桶、B2的$4K$块、B3的$8K$块
\FOR{每个起点}
\STATE 保留原方案；各尝试迁块、同核合并、调序及容量风险块的前后移动
\STATE 检查覆盖和依赖，规范化并按计划内容去重
\ENDFOR
\STATE 按$(\widehat F_B,\widehat Q_B,\Pi)$排列候选
\STATE 将整图、分量装桶、上一核继承计划和问题一热启动计划置于去重后的必评队列
\STATE 令$m$为必评数，完整调用额度为$\max\{\min(q,2),m+1\}$
\WHILE{尚有候选且未达到常规完整调用额度}
\STATE 对当前候选做展开检查；若失败，记录错误并继续
\STATE 运行B场景事件模拟，保存成功结果或失败文本
\ENDWHILE
\STATE 若当前最优慢于上一核继承计划，补评尚未评价的回退计划
\STATE 在成功集合按$(F_B,D_B)$字典序选$X_B^*$
\end{algorithmic}\end{algorithm}

超过20000个计算操作的五张图也沿用B0--B3构造并纳入100图结果。
\FloatBarrier

\subsection{问题二结果分析}
\subsubsection{逐核加速比}
加速比仍以整图单核A0的$T_{1,i}$为基准，计算$T_{1,i}/F_B(X^{\rm inh}_{B,iK})$；单核曲线点按定义显示为1。表\ref{tab:b-results}中各核数都没有慢于基准的图。图\ref{fig:speedup-b}的浅色带为100张图重采样1500次得到的均值95\%区间。
\begin{table}[H]\centering\small\caption{问题二：100图逐核加速比（继承方案）}\label{tab:b-results}
\begin{tabular}{crrrr}\toprule
核数 & 平均加速比 & 中位数 & 慢于A0的图数 & 首个成功至最终平均降幅\\\midrule
1 & 1.0000 & 1.0000 & 0 & 1.6528\%\\
2 & 1.7122 & 1.9647 & 0 & 0.3034\%\\
3 & 2.3140 & 2.8340 & 0 & 1.5168\%\\
4 & 2.9183 & 3.6353 & 0 & 2.6911\%\\
5 & 3.3862 & 4.0046 & 0 & 6.3530\%\\\bottomrule
\end{tabular}\end{table}
\begin{figure}[!htbp]\centering
\includegraphics[width=0.91\linewidth,height=0.43\textheight,keepaspectratio]{figures/speedup_B.pdf}
\caption{问题二逐核加速比：均值、中位数、均值95\%重采样区间及4至5核重点区间}\label{fig:speedup-b}
\end{figure}

五核平均加速比为3.3862，中位数4.0046，均值的95\%重采样区间为$[3.0686,3.6824]$。由四核到五核，均值从2.9183升至3.3862，增加0.4679；65张图工期缩短、35张图保持不变，没有图变慢。五核有21张图的最终候选缩短了工期。
五核继承方案的平均额外搬运为15596420.30 B，其中结构附加4814112.38 B、Spill 10782307.92 B；38张图出现非零Spill，均值主要由少数大图的多次恢复拉动。

两场景五核继承方案的平均加速比分别为3.4132和3.3862，相差0.0270。逐图比较，B比A快22图、持平68图、慢10图。B更慢的图上，核级驻留节省的读入被Spill、跨核同步或关键路径上的等待抵消。图\ref{fig:b-ab-tradeoff}用固定五核方案逐图配对，横轴为额外搬运差，纵轴为B/A工期比；大多数图的额外搬运差和工期比不成单调关系，说明核级驻留的收益主要由核内复用决定，字节数高低并不直接决定工期快慢。
\samplefigure{ab_tradeoff_scatter}{五核固定核数方案的逐图A/B对照：额外搬运差与工期比}{fig:b-ab-tradeoff}

\begin{table}[H]\centering\small\caption{五核代表图在A/B场景的工期及额外搬运分解（固定核数方案）}\label{tab:b-cases}
\begin{tabular}{llrrr}\toprule
图 & 场景及候选 & 工期/cycle & 结构附加/B & Spill/B\\\midrule
\Case{055} & A：A2迁块 & 110002 & 199780 & 0\\
 & B：B1原方案 & 25635 & 231424 & 0\\
\Case{053} & A：A1原方案 & 492699 & 3104256 & 3814656\\
 & B：B0整图 & 1472686 & 0 & 5405536\\\bottomrule
\end{tabular}\end{table}

图\ref{fig:b-case055-timeline}展开\Case{055}五核B方案的事件时间线。
\samplefigure{b_case055_official_timeline}{\Case{055}五核B方案的核级Task、COPY及在途搬运时间线}{fig:b-case055-timeline}

该图B工期仅为A的约0.233倍，结构附加搬运却从199780 B增加到231424 B，两者均无Spill：B更快并非因为字节更少，而是Task组织规则改变后，决定完成时间的链路缩短了。\Case{053}方向相反：固定五核时B选中整图B0，结构附加为零，但Spill达到5405536 B，工期约为A的2.99倍。零结构通信抵不过失去多核并行与长驻留带来的恢复代价；继承后该图B的五核结果改用二核方案，工期828335 cycle。

\subsubsection{容量恢复算例}
结构附加搬运相同时，可以单独观察块结构与核内次序对Spill的影响。\Case{008}单核B3起点与其相邻块合并变体的结构附加均为0，前者Spill为10562688 B、工期500608 cycle，后者Spill为7059072 B、工期405761 cycle。合并改变了块结构及核内优先次序，Spill减少3503616 B，工期缩短94847 cycle。
\begin{table}[H]\centering\caption{\Case{008}单核问题二两份成功候选}\label{tab:b-case008}
\begin{tabular}{lrrr}\toprule
候选 & 工期/cycle & 结构附加/B & Spill/B\\\midrule
B3原方案 & 500608 & 0 & 10562688\\
B3合并变体 & 405761 & 0 & 7059072\\\bottomrule
\end{tabular}\end{table}

\subsection{问题二结论}
问题二的核级接收域把同核消费归并为一份结构读入，式\eqref{eq:b-bytes}给出跨核COPY的计量，式\eqref{eq:b-physical-live}和式\eqref{eq:b-spill}说明容量恢复如何改变工期。按继承方案，五核平均加速比为3.3862，所有图均不慢于整图单核基准。

提交计划包含操作归属与核心列表。问题三的同计划配对以同核数下问题二的固定核数方案$X_B^*$为起点（未经继承），在相同计划上加入只读Cache，再比较硬件作用与方案适配。各类结果记录见附录表\ref{tab:b-artifacts}。

\subsection{问题二灵敏度分析}
问题二的片上驻留对分块方式敏感。表\ref{tab:b-case008}的两份候选在相同硬件和核数下只差一次相邻块合并，Spill与工期同时明显变化，说明核级Task内物理副本的生存期随块边界和操作次序改变。增加核心数的影响见表\ref{tab:b-results}和图\ref{fig:speedup-b}：继承方案下四至五核有65图工期缩短、35图不变。
